\chapter{Conclusiones}

Este trabajo completó la construcción y validación empírica del instrumento de caracterización CPU--GPU de la Fase 1; el conjunto DVFS de CPU quedó cerrado y el de GPU aún requiere su recorrido experimental completo. Las dos campañas CPU reunieron 1470 corridas aceptadas y 1\,309\,755 intervalos \textit{uncore} elegibles con FLOPs medidos por hardware, bytes reales de DRAM y frecuencia observada. Con el turbo desactivado, los nueve niveles fijos son físicamente distintos (de 3.2 a 0.8~GHz) y cada ventana se acepta por su frecuencia observada, de modo que estos datos sirven tanto para entrenar la clasificación como para comparar tiempo, energía y EDP entre niveles. En GPU, el catálogo quedó caracterizado y auditado por precisión, y \texttt{-lgc} retiene el reloj bajo carga real, aunque falta aceptar la campaña multi-frecuencia completa.

La Fase 2 produjo el clasificador de fase CPU y su evaluación por familia algorítmica retirada. El modelo final es un XGBoost con doce tasas derivadas de contadores de rendimiento y la frecuencia observada, ajustado con peso total igual por celda de familia y clase, y con una salida \textit{revisar} cuando su confianza no alcanza 0.90. Ante una familia algorítmica que no intervino en ninguna etapa del ajuste clasifica correctamente el 72.7\% de los intervalos en la métrica balanceada por celda, frente al 46.4\% de la referencia sin capacidad de discriminar, y alcanza 75.5\% si se le permite abstenerse en el 33.4\% de los casos. La comparación controlada mostró que ni el algoritmo de aprendizaje, ni el número de variables de entrada, ni la búsqueda de hiperparámetros con Optuna modifican ese resultado más allá del intervalo de confianza, mientras que la diversidad del catálogo sí lo hace, con una ganancia estimada de 0.039 por duplicación del número de familias. La tabla de política de frecuencia derivada de las dos campañas resultó \textit{no actuar} en ambas clases: la potencia de piso, de unos 85~W y el 83\% de la potencia a 3.2~GHz, hace que bajar el reloj solo compense en cargas con exponente de tiempo $\alpha\lesssim 0.22$, que cumplen cinco kernels tipo STREAM con ganancias de 4 a 12\%, y el mejor caso posible con un clasificador perfecto sería de 2.5\% en promedio. El reanálisis retrospectivo de campañas GPU con NVML obtuvo F1 macro externo de 0.663 sobre cuatro familias no vistas, frente a 0.426 de la regla mayoritaria. La implementación completa del agente y la evaluación de EDP corresponden a las Fases 3 y 4.

El aporte metodológico central de esta fase, más allá de los datos concretos recolectados, es haber demostrado, y no solo enunciado, que ningún mecanismo de este instrumento se dio por funcional sin verificación directa contra el estado real del hardware: ni la disponibilidad de un contador de PMU, ni el presupuesto físico de contadores simultáneos, ni su persistencia en el tiempo frente a mecanismos ajenos al instrumento, ni la escritura de frecuencia de CPU, ni la precisión aritmética declarada de un kernel de GPU, ni la calibración de los propios techos del modelo Roofline. En varias ocasiones documentadas en este trabajo, esa verificación expuso una falla real que una confianza no verificada en la documentación del fabricante, en una medición previa ya dada por válida, o en el diseño original del instrumento, habría dejado pasar silenciosamente hacia el conjunto de datos.

\section{Limitaciones}

Este trabajo reconoce las siguientes limitaciones. Primera, la verdad CPU tiene la granularidad del intervalo \textit{uncore}, de aproximadamente 10~ms o más, y no la de las ventanas de núcleo de aproximadamente 1~ms. Segunda, los 1.31 millones de intervalos proceden de 30 familias y no constituyen unidades independientes de generalización. El límite de 1000 intervalos por celda familia--clase puede omitir submodos minoritarios dentro de una celda, aunque preserve todas las familias y clases observadas. Tercera, la abstención mejora la calidad agregada dentro de la cobertura, pero no elimina errores de alta confianza en familias no vistas. Cuarta, el presupuesto simultáneo de PMU depende de mecanismos del sistema operativo y requiere repetir su verificación si cambia el entorno. Quinta, el control de frecuencia CPU por núcleo exige desactivar el turbo global, lo que depende de un permiso específico del nodo. Sexta, el clasificador GPU parte de 394 corridas y 13 familias, con predominio \textit{memory-bound}; su validación externa contiene solo cuatro familias inéditas. Séptima, NVML no mide directamente los contadores microarquitectónicos que definen el régimen Roofline y las variables temporales no transfirieron entre campañas. Octava, todavía no existe una campaña GPU multi-frecuencia completa aceptada.

\section{Trabajo Futuro}

En el corto plazo, tres tareas siguen directamente habilitadas. Primera, en CPU, evaluar el modelo congelado sobre familias compute inéditas reservadas como prueba externa y llevar la tabla de política (hoy \textit{no actuar} en ambas clases) a las aplicaciones compuestas de la Fase 3. Segunda, desbloquear la campaña GPU y medir la intrusión de CUPTI Activity antes de utilizar sus trazas para alinear etiquetas por fase o ventana. El perfilado de CUPTI debe conservarse para construir y auditar etiquetas, mientras que NVML debe permanecer como entrada del clasificador ligero. Tercera, ejecutar una nueva campaña GPU con familias inéditas y evaluar una única vez el modelo congelado en esa campaña, sin reajustar umbrales ni seleccionar características con sus resultados.
